\documentclass[11pt, a4paper, openany]{article}
\usepackage{ctex}
\usepackage{amsmath, amssymb, amsthm}
\usepackage{geometry}
\geometry{margin=1.5in}

% 定理环境定义
\newtheorem{definition}{定义}[section]
\newtheorem{theorem}{定理}[section]
\newtheorem{corollary}{推论}[section]
\newtheorem{lemma}{引理}[section]
\newtheorem{example}{例题}[section]
\newtheorem{remark}{注}[section]
\newtheorem{errorpoint}{易错点}[section]

% 取消多余缩进与间距
\setlength{\parindent}{0em}
\setlength{\parskip}{0.5em}
\renewcommand{\baselinestretch}{1.2}

\begin{document}

\title{数学分析（华东师大第五版）复习笔记}
\maketitle

\section{实数集与函数}
\subsection{实数}
\begin{definition}
\textbf{实数}是有理数和无理数的总称，其中有理数可表示为分数$\frac{p}{q}$（$p,q$为整数，$q>0$且互质），也可表示为有限十进小数或无限十进循环小数；无理数是无限十进不循环小数。
\end{definition}

\begin{remark}
为统一表示，约定：将非零有限小数$x=a_0.a_1a_2\cdots a_n$（$a_n\neq0$）记为$a_0.a_1a_2\cdots(a_n-1)999\cdots$，将正整数$m$记为$(m-1).999\cdots$，负有限小数则先对其绝对值做上述转换后再加负号，该表示称为实数的\textbf{正规表示}，保证每个实数对应唯一的无限小数表示。
\end{remark}

\begin{definition}
两个实数$x=a_0.a_1a_2\cdots$，$y=b_0.b_1b_2\cdots$（均为正规表示），若对所有非负整数$k$都有$a_k=b_k$，则称$x=y$；若存在非负整数$n$，使得$a_k=b_k$对所有$k<n$成立且$a_n>b_n$，则称$x$大于$y$，记为$x>y$。
\end{definition}

\begin{theorem}[实数的稠密性]
任意两个不相等的实数之间必有有理数，也必有无理数。
\end{theorem}

\begin{definition}
实数$a$的\textbf{绝对值}定义为：
$$|a|=\begin{cases}
a, & a\geq0, \\
-a, & a<0.
\end{cases}$$
其几何意义是数轴上点$a$到原点的距离。
\end{definition}

\begin{theorem}[绝对值的性质]
对任意实数$a,b$，有：
\begin{enumerate}
    \item $|a|\geq0$，当且仅当$a=0$时$|a|=0$；
    \item $|-a|=|a|$；
    \item $-|a|\leq a\leq|a|$；
    \item \textbf{三角不等式}：$|a|-|b|\leq|a\pm b|\leq|a|+|b|$；
    \item $|ab|=|a||b|$；
    \item 若$b\neq0$，则$\left|\frac{a}{b}\right|=\frac{|a|}{|b|}$。
\end{enumerate}
\end{theorem}

\begin{example}
证明对任意实数$a,b$，有$||a|-|b||\leq|a-b|$。
\begin{proof}
由三角不等式$|a|=|(a-b)+b|\leq|a-b|+|b|$，移项得$|a|-|b|\leq|a-b|$；同理$|b|=|(b-a)+a|\leq|b-a|+|a|=|a-b|+|a|$，移项得$|b|-|a|\leq|a-b|$，即$|a|-|b|\geq-|a-b|$。结合两式得$||a|-|b||\leq|a-b|$。
\end{proof}
\end{example}

\subsection{实数集的界与确界原理}
\begin{definition}
设$S$为$\mathbb{R}$中的非空数集，若存在实数$M$，使得对所有$x\in S$都有$x\leq M$，则称$M$是$S$的一个\textbf{上界}，$S$称为\textbf{有上界的数集}；若存在实数$L$，使得对所有$x\in S$都有$x\geq L$，则称$L$是$S$的一个\textbf{下界}，$S$称为\textbf{有下界的数集}。若$S$既有上界又有下界，则称$S$是\textbf{有界集}，否则称$S$为\textbf{无界集}。
\end{definition}

\begin{remark}
数集的上界、下界不唯一：若$M$是$S$的上界，则所有大于$M$的实数都是$S$的上界；同理若$L$是$S$的下界，则所有小于$L$的实数都是$S$的下界。
\end{remark}

\begin{definition}
设$S$为$\mathbb{R}$中的非空有上界数集，若实数$\eta$满足：
\begin{enumerate}
    \item 对所有$x\in S$，有$x\leq\eta$（即$\eta$是$S$的上界）；
    \item 对任意$\varepsilon>0$，存在$x_0\in S$使得$x_0>\eta-\varepsilon$（即任何小于$\eta$的数都不是$S$的上界），
\end{enumerate}
则称$\eta$是$S$的\textbf{上确界}，记为$\eta=\sup S$。
\end{definition}

\begin{definition}
设$S$为$\mathbb{R}$中的非空有下界数集，若实数$\xi$满足：
\begin{enumerate}
    \item 对所有$x\in S$，有$x\geq\xi$（即$\xi$是$S$的下界）；
    \item 对任意$\varepsilon>0$，存在$x_0\in S$使得$x_0<\xi+\varepsilon$（即任何大于$\xi$的数都不是$S$的下界），
\end{enumerate}
则称$\xi$是$S$的\textbf{下确界}，记为$\xi=\inf S$。
\end{definition}

\begin{errorpoint}
数集的上确界、下确界不一定属于该数集，例如开区间$(0,1)$的上确界是$1$，下确界是$0$，二者都不属于$(0,1)$；若上确界属于数集，则称其为$S$的最大值，同理若下确界属于数集，则称其为$S$的最小值。
\end{errorpoint}

\begin{theorem}[确界原理]
非空有上界的数集必有唯一的上确界；非空有下界的数集必有唯一的下确界。
\end{theorem}

\begin{remark}
确界原理是实数完备性的基本定理之一，仅在实数域成立，有理数域不满足该原理，例如所有平方小于2的正有理数构成的集合在有理数域内有上界（如2），但不存在有理数上确界。
\end{remark}

\begin{example}
求数集$S=\left\{y\left|y=2-x^2, x\in\mathbb{R}\right.\right\}$的上、下确界。
\begin{proof}
对任意$x\in\mathbb{R}$，$x^2\geq0$，故$y=2-x^2\leq2$，即2是$S$的上界；对任意$\varepsilon>0$，取$x_0=\sqrt{\varepsilon/2}$，则$y_0=2-(\varepsilon/2)=2-\varepsilon/2>2-\varepsilon$，满足上确界定义，故$\sup S=2$。
对任意实数$M>0$，取$x_0=\sqrt{M+3}$，则$y_0=2-(M+3)=-M-1<-M$，说明$S$无下界，因此下确界不存在（或记为$\inf S=-\infty$）。
\end{proof}
\end{example}

\subsection{函数概念}
\begin{definition}
设$D$是$\mathbb{R}$中的非空子集，若存在对应法则$f$，使得对每个$x\in D$，都有唯一确定的实数$y$与之对应，则称$f$是定义在$D$上的\textbf{函数}，记为$y=f(x), x\in D$，其中$D$称为函数的\textbf{定义域}，$x$称为自变量，$y$称为因变量，所有函数值构成的集合$f(D)=\left\{y\left|y=f(x),x\in D\right.\right\}$称为函数的\textbf{值域}。
\end{definition}

\begin{remark}
函数的两个核心要素是\textbf{定义域}和\textbf{对应法则}，当且仅当两个函数的定义域和对应法则完全相同时，二者是同一个函数，与自变量、因变量的符号无关。
\end{remark}

\begin{definition}
设函数$y=f(u)$的定义域为$D_f$，函数$u=g(x)$的定义域为$D_g$，且$g(D_g)\subset D_f$，则称函数
$$y=f(g(x)),\quad x\in D_g$$
为$f$与$g$的\textbf{复合函数}，其中$u$称为中间变量，记为$f\circ g$。
\end{definition}

\begin{definition}
设函数$y=f(x)$的定义域为$D$，值域为$f(D)$，若对每个$y\in f(D)$，都存在唯一的$x\in D$使得$f(x)=y$，则按此对应法则得到的定义在$f(D)$上的函数称为$f$的\textbf{反函数}，记为$x=f^{-1}(y), y\in f(D)$，通常改写为$y=f^{-1}(x), x\in f(D)$。
\end{definition}

\begin{theorem}
严格单调函数必有反函数，且反函数与原函数单调性相同。
\end{theorem}

\subsection{具有某些特性的函数}
\begin{definition}
设$f$为定义在$D$上的函数，若存在正数$M$，使得对所有$x\in D$都有$|f(x)|\leq M$，则称$f$为$D$上的\textbf{有界函数}；否则称$f$为$D$上的\textbf{无界函数}。若存在实数$M$，使得对所有$x\in D$都有$f(x)\leq M$，则称$f$在$D$上有上界；若存在实数$L$，使得对所有$x\in D$都有$f(x)\geq L$，则称$f$在$D$上有下界。
\end{definition}

\begin{remark}
函数$f$在$D$上有界的充要条件是$f$在$D$上既有上界又有下界。
\end{remark}

\begin{definition}
设$f$为定义在$D$上的函数，若对任意$x_1,x_2\in D$，当$x_1<x_2$时总有$f(x_1)\leq f(x_2)$，则称$f$为$D$上的\textbf{单调递增函数}；若总有$f(x_1)<f(x_2)$，则称$f$为$D$上的\textbf{严格单调递增函数}。类似可定义单调递减函数和严格单调递减函数，单调递增和单调递减函数统称为\textbf{单调函数}。
\end{definition}

\begin{definition}
设$D$为关于原点对称的数集，若对所有$x\in D$都有$f(-x)=f(x)$，则称$f$为\textbf{偶函数}；若对所有$x\in D$都有$f(-x)=-f(x)$，则称$f$为\textbf{奇函数}。
\end{definition}

\begin{remark}
偶函数的图像关于$y$轴对称，奇函数的图像关于原点对称。
\end{remark}

\begin{definition}
设$f$为定义在$D$上的函数，若存在非零常数$T$，使得对所有$x\in D$有$x\pm T\in D$且$f(x+T)=f(x)$，则称$f$为\textbf{周期函数}，$T$为$f$的一个周期；若周期中存在最小的正数，则称其为$f$的\textbf{基本周期}，通常简称周期。
\end{definition}

\begin{errorpoint}
并非所有周期函数都有基本周期，例如常值函数$f(x)=C$是周期函数，任意非零实数都是其周期，但不存在最小正周期；狄利克雷函数也是周期函数，任意正有理数都是其周期，同样无基本周期。
\end{errorpoint}

\end{document}

















